\newpage
\section{Introduction}
\label{sec:introduction}

The emission of EM radiation from a body can be (in which cases?) attributed to photonic and electronic contributions separately [Novotny], namely

\begin{align}
I(\hbar\omega) &= I_{\text{ph}}(\hbar\omega)\cdot I_{\text{e}}(\hbar\omega)
\label{eq:intro_emission_split}
\end{align}

In this study we focus on the latter. Electronic transitions which begin at a high energy state and end up in a lower energy state result in the spontaneous emission of photons with \textbf{appropriate energy}. Transitions such as these can occur either inside the conduction band (assisted by phonons) or between bands - intraband and interband transitions respectively\footnote{Technically transitions can occur inside the valence band, but these are very unlikely and their contribution at standard cases is negligible.}.

\begin{equation}
I_{\text{e}} = I_{\text{cc}} + I_{\text{cv}}
\label{eq:intro_emission_components}
\end{equation}

Out materials of interest are metals (specifically gold) in which the predominant electronic contribution to emission comes from intraband transitions\footnote{TODO: add footnote text.}. Therefore, in a state of thermal equilibrium we also don't expect measurable interband presence in emission spectrum from metals [Marini?].

\begin{tcolorbox}[title={Note: Main Thesis}]
In non-equilibrium scenarios, however, such as in the steady-state case of a metal illuminated by monochromatic radiation, incoming photons of sufficient energy might be able to generate a significant hole population in the valence band which would allow for interband transitions to be more prominent in the emission spectrum of the metal.
\end{tcolorbox}

Non-equilibrium electronic distribution for the above scenario has been derived and applied to the case of intraband transitions [Sivan \& Dubi]. These calculations, however, enjoyed the approximate isotropy of the conduction band which allows one to transition to the more lenient energy space.

The valence band is generally anisotropic and when performing interband calculations one has to remain in k-space, which is what we are going to do.

Once we have established that our calculations in momentum space align with those performed in energy space - both in equilibrium and non-equilibrium settings - we proceed to perform the interband scenario. In the end, hopefully, we will be able to provide a more complete emission spectrum to both cases.

\subsection{Gold Nanoparticles}
We focus on interband transitions from the X and L points in k-space whose approximate symmetry can be used to simplify calculations. They also account for lion-share of the interband transitions, but not all of them.

\begin{figure}[H]
  \centering
  \includegraphics[width=0.35\textwidth]{pasted_image_20260121090737.png}
  \caption{Pasted image 20260121090737.}
  \label{fig:gold_nanoparticles}
\end{figure}

The overall electronic contribution to emission is therefore given by

\begin{equation}
\Gamma = \Gamma_{\text{cc}} + 6\Gamma_{\text{cv}}^{\text{X}} + 8\Gamma_{\text{cv}}^{\text{L}}
\label{eq:intro_total_emission}
\end{equation}

\subsection{Progress Report, Tasks \& Questions}

\subsubsection{Progress Report}
\begin{itemize}
  \item Settled on a convergent numerical scheme for the calculation of the emission integral. (\hyperref[sec:numeric_convergence]{\ref*{sec:numeric_convergence}~\nameref*{sec:numeric_convergence}})
  \item Established approximate nature of the analytic expressions for intraband emission (thermal and non-equilibrium). (\hyperref[sec:analytic_approximations]{\ref*{sec:analytic_approximations}~\nameref*{sec:analytic_approximations}})
  \begin{itemize}
    \item Outlined its region of validity.
    \item Found an exact analytic expression for thermal intraband with constant eDOS.
  \end{itemize}
  \item Quantified the discrepancy between constant eDOS and energy-dependent eDOS. (\hyperref[fig:constant_edos_approximation]{Figure~\ref*{fig:constant_edos_approximation}})
  \item Recognized the difficulty and inefficiency of working in k-space and having to approximate the delta function numerically (\hyperref[app:numeric_delta_approximation]{\ref*{app:numeric_delta_approximation}~\nameref*{app:numeric_delta_approximation}})
  \item Found the proper transformations that take us to energy space, even with non-isotropic dispersion relations (\hyperref[app:quadratic_intraband_transitions]{\ref*{app:quadratic_intraband_transitions}~\nameref*{app:quadratic_intraband_transitions}}, \hyperref[app:quadratic_interband_transitions]{\ref*{app:quadratic_interband_transitions}~\nameref*{app:quadratic_interband_transitions}})
\end{itemize}

\subsubsection{Tasks}
\begin{itemize}
  \item[$\boxtimes$] Improve abstract
  \item[$\boxtimes$] add a more quantitative view of the constant eDOS mismatch
  \item[$\square$] Focus on lower emission energies.
  \item[$\square$] Reproduce Rossei's absorption calculation
  \item[$\square$] Calculate intraband emission using Rossei's bands and compare the free electron approximation (parabolic band)
  \item[$\square$] Calculate interband emission using Rossei's bands and compare to absorption - in thermal eq emission = absorption\* plank (Kirchhoff).
  \item[$\square$] ?
\end{itemize}

\subsubsection{Questions}
\begin{itemize}
  \item Absorption and emission happen only at skin-depth. The statistical transition from a single emitter to continuum includes counting emitters in the particle's entire volume. Isn't it more appropriate to calculate specific emission?
  \item We're either calculating electronic transitions in bulk metals in which case we have to account for phonon-electron interaction for momentum conservation to hold.
  \item Or we're calculating electronic transitions in nano-particles in which a first order transition is opened due to the confinement of the wave-function in space and the spreading of momentum. (Landau damping)
  \item In testing theoretical predictions against experimental evidence, shouldn't the photonic DOS 
  \item In experimental settings I imagine LDOPS (which are generally not that of vacuum?) affect the emission spectrum both qualitatively and quantitatively - how do we address this when drawing comparisons?
\end{itemize}
